%! Unit 1: Asking Questions with Data: Study Design — Quiz 1 — Sample Quiz (Lessons 1.0–1.3) — blank
% Standalone assessment on the unit-test model (LESSON_SHAPE.md section 6): 10pt,
% \parthead{Part …} strips, Part A Vocabulary · B Multiple Choice · C Short Answer &
% Computation · D Extended Response, \workrowsep 5pt identical in blank and key, and
% \namedateperiod kept because it is taken in a testing setting. Merged into no packet.
%
% BLUEPRINT — 19 items, 50 points, Lessons 1.0–1.3 only:
%   A vocabulary matching  8 items ×1 =  8   (1.0 ×3, 1.1 ×1, 1.2 ×2, 1.3 ×2)
%   B multiple choice      6 items ×2 = 12   (1.0, 1.1, 1.2 ×2, 1.3 ×2)
%   C short answer         4 items ×5 = 20   (C1 1.0/1.1 · C2 1.1 · C3 1.2 · C4 1.3)
%   D extended response    1 item ×10 = 10   (1.3 crux + 1.2 parameter/statistic)
% Parallel form to the real quiz: same blueprint and ideas, different numbers and contexts,
% reshuffled vocabulary letters. Every number verified in Python before authoring.
% Contexts are deliberately NEW (Millbrook, Stonebridge) — Riverbend, Lakeside, Harbor Point,
% Northgate, Cedar Ridge, and Bayside all appear in lessons 1.0–1.3 or the unit test.
%
% THE KEY IS GENERATED, NEVER HAND-EDITED:
%   python3 templates/lesson/mkkey.py unit01/quiz01/sample_quiz/main.tex \
%       unit01/quiz01/sample_quiz_key/main.tex unit01/quiz01/sample_quiz_key/answers.json
\documentclass[10pt]{article}
\usepackage{saar-article}
\usepackage{saar-boxes}

% Work blocks on a quiz need handwriting room, not typeset spacing. This moves the blank and
% the key together, so it cannot open a page mismatch.
\setlength{\workrowsep}{5pt}

% Part-divider strip
\newcommand{\parthead}[1]{%
  \vspace{6pt}\noindent
  \begin{headlinebox}{cerulean}{\color{white}\bfseries #1}\end{headlinebox}%
  \vspace{4pt}}

\begin{document}

\pageheader{Unit 1: Asking Questions with Data: Study Design}{Quiz 1 --- Sample Quiz --- Lessons 1.0--1.3}
\namedateperiod

\vspace{0.10in}
\begin{remindbox}
\small
\textbf{This is a sample quiz.} It mirrors the real Quiz 1 in length, format, and the ideas it
covers, but the numbers and contexts differ. The quiz covers \textbf{Lessons 1.0--1.3 only}.
Work every problem with no notes, then check your answers against the key.
\textbf{50 points:} Part A $8$, Part B $12$, Part C $20$, Part D $10$.
\end{remindbox}

% ======================================================================
\parthead{Part A --- Vocabulary \hfill 8 questions, 1 point each}

\small
Write the letter of the matching description next to each term. Each letter is used once.

\vspace{2pt}
\begin{minipage}[t]{0.36\linewidth}
\begin{enumerate}[leftmargin=1.9em, itemsep=3.5pt, topsep=1pt]
  \item \blank{0.8cm} Individual
  \item \blank{0.8cm} Categorical variable
  \item \blank{0.8cm} Statistical question
  \item \blank{0.8cm} Population
  \item \blank{0.8cm} Parameter
  \item \blank{0.8cm} Statistic
  \item \blank{0.8cm} Sampling frame
  \item \blank{0.8cm} Cluster sample
\end{enumerate}
\end{minipage}\hfill
\begin{minipage}[t]{0.61\linewidth}
\begin{enumerate}[label=\Alph*., leftmargin=1.9em, itemsep=3.5pt, topsep=1pt]
  \item A question whose answers are expected to vary from individual to individual.
  \item A number that describes the population. Usually unknown.
  \item The list of individuals a sample can actually be drawn from.
  \item One person or object that data are collected about --- one row of a data table.
  \item A few whole groups are chosen at random, and everyone in them is studied.
  \item Every individual the study wants to describe.
  \item A variable that records a label or group name, summarized with counts and percents.
  \item A number computed from the sample. It changes from sample to sample.
\end{enumerate}
\end{minipage}
\normalsize

% ======================================================================
\vspace{4pt}
\boxguard[8]
\parthead{Part B --- Multiple Choice \hfill 6 questions, 2 points each}

\small
\begin{enumerate}[leftmargin=2.2em, itemsep=5pt, topsep=2pt, start=9]

\item Millbrook High School records each student's \textbf{advisory group number}, from $1$ to
$30$. This variable is
\begin{enumerate}[label=(\Alph*), leftmargin=1.8em, itemsep=1pt, topsep=1pt]
  \item quantitative, because the values are numbers.
  \item categorical, because the numbers are labels for groups.
  \item quantitative, because the groups can be put in order.
  \item neither, because a group number is not really data.
\end{enumerate}
\hfill \textbf{Answer:} \blank{0.9cm}

\item Which of these is a \textbf{statistical question}?
\begin{enumerate}[label=(\Alph*), leftmargin=1.8em, itemsep=1pt, topsep=1pt]
  \item How many minutes did I spend on homework last night?
  \item How many minutes do Millbrook students spend on homework on a school night?
  \item How many seats does the Millbrook gym hold?
  \item Does Millbrook start classes at 8 a.m.?
\end{enumerate}
\hfill \textbf{Answer:} \blank{0.9cm}

\item Stonebridge Recreation Center's own membership records show that \textbf{exactly $45\%$
of all $1{,}200$ members} renew early. That $45\%$ is a
\begin{enumerate}[label=(\Alph*), leftmargin=1.8em, itemsep=1pt, topsep=1pt]
  \item statistic, because it is written as a percent.
  \item statistic, because $1{,}200$ is a large number.
  \item parameter, because it describes every member of the population.
  \item parameter, because the records are official.
\end{enumerate}
\hfill \textbf{Answer:} \blank{0.9cm}

\item Two different random samples of $80$ Stonebridge members give $35\%$ and $41\%$ arriving
before 7 a.m. This result shows
\begin{enumerate}[label=(\Alph*), leftmargin=1.8em, itemsep=1pt, topsep=1pt]
  \item that one of the two samples was collected incorrectly.
  \item that the parameter changed between the two surveys.
  \item sampling variability --- the parameter did not move, and both are honest estimates.
  \item that the true value must be exactly halfway between them.
\end{enumerate}
\hfill \textbf{Answer:} \blank{0.9cm}

\item A teacher surveys her own second-period class because it is the easiest group to reach.
Her sample is
\begin{enumerate}[label=(\Alph*), leftmargin=1.8em, itemsep=1pt, topsep=1pt]
  \item a simple random sample, because she picked no favorites.
  \item a systematic sample, because her class list is in order.
  \item a cluster sample, because she surveyed one whole class.
  \item a convenience sample, because a student in another period could never be chosen.
\end{enumerate}
\hfill \textbf{Answer:} \blank{0.9cm}

\item Which sentence states the difference between a \textbf{stratum} and a \textbf{cluster}?
\begin{enumerate}[label=(\Alph*), leftmargin=1.8em, itemsep=1pt, topsep=1pt]
  \item A stratum is alike inside and differs from other strata; a cluster is meant to be a
        small mix of the whole.
  \item A stratum is a small mix of the whole; a cluster is alike inside.
  \item A stratum is chosen by chance; a cluster is chosen by convenience.
  \item There is no difference --- both mean a group inside the population.
\end{enumerate}
\hfill \textbf{Answer:} \blank{0.9cm}

\end{enumerate}
\normalsize

% ======================================================================
\newpage
\parthead{Part C --- Short Answer \& Computation \hfill 4 questions, 5 points each}

\small
\begin{enumerate}[leftmargin=2.2em, itemsep=9pt, topsep=2pt, start=15]

% ---------------------------------------------------------------- C1 (1.0 / 1.1)
\item \textbf{Millbrook High School} has $720$ students. For each student the office records
advisory group number, minutes on the bus this morning, whether the student buys lunch, and
preferred spirit-week theme.

\vspace{2pt}
(a) Who are the \textbf{individuals} in this data set? \blank{6.2cm}

\vspace{3pt}
(b) Classify each variable.

\vspace{2pt}
\renewcommand{\arraystretch}{1.25}
\begin{tabularx}{\linewidth}{@{} X c @{}}
\rowcolor{frost}
\textbf{Variable} & \textbf{Categorical or quantitative?} \\
\midrule
Advisory group number & \blank{4.4cm} \\
Minutes on the bus this morning & \blank{4.4cm} \\
Buys lunch (yes / no) & \blank{4.4cm} \\
Preferred spirit-week theme & \blank{4.4cm} \\
\end{tabularx}

\vspace{4pt}
(c) The office's first draft question is \emph{``How long is the bus ride?''} Explain why that
is not a statistical question, then rewrite it so that it is.

\par\writeline
\par\writeline

% ---------------------------------------------------------------- C2 (1.1)
\boxguard[20]
\item \textbf{Millbrook's student council} asks $60$ students to name a preferred spirit-week
theme. The counts are below.

\vspace{2pt}
(a) Complete the relative frequency column as a percent.

\vspace{2pt}
\renewcommand{\arraystretch}{1.25}
\begin{tabularx}{\linewidth}{@{} X c c @{}}
\rowcolor{frost}
\textbf{Theme} & \textbf{Frequency} & \textbf{Relative frequency} \\
\midrule
Decades     & $24$ & \blank{2.6cm} \\
Color wars  & $18$ & \blank{2.6cm} \\
Sports day  & $12$ & \blank{2.6cm} \\
Movie night & $6$  & \blank{2.6cm} \\
\end{tabularx}

\vspace{4pt}
(b) Write a \textbf{conclusion in context} about the most popular theme, naming the group it
applies to.

\par\writeline
\par\writeline

\vspace{2pt}
(c) The council writes: \emph{``Most Millbrook students want a decades theme.''} Give the two
things that are wrong with that sentence.

\par\writeline
\par\writeline

% ---------------------------------------------------------------- C3 (1.2)
\boxguard[18]
\item \textbf{Stonebridge Recreation Center} has $1{,}200$ members. The manager surveys $80$
members chosen at random; $28$ of them arrive before 7 a.m.

\vspace{2pt}
(a) Population: \blank{4.4cm} \quad Sample size: \blank{1.8cm}

\vspace{3pt}
Sample: \blank{9.0cm}

\vspace{3pt}
(b) What percent of the sample arrives before 7 a.m., and about how many of all $1{,}200$
members does that estimate?

\begin{work}
  \dfrac{28}{80} &= 0.35 = 35\% \\
  0.35 \times 1200 &= 420 \text{ members}
\end{work}

(c) Is the $35\%$ a parameter or a statistic? Say how you know.

\par\writeline

(d) Name one \textbf{constraint} that keeps the manager from taking a census of all $1{,}200$
members.

\par\writeline

% ---------------------------------------------------------------- C4 (1.3)
\newpage
\item \textbf{Millbrook High School} has $720$ students, a numbered list of all of them, and
$30$ advisory groups of $24$ students each.

\vspace{2pt}
(a) Name the sampling technique used in each row.

\vspace{2pt}
\renewcommand{\arraystretch}{1.3}
\begin{tabularx}{\linewidth}{@{} X c @{}}
\rowcolor{frost}
\textbf{What the researcher did} & \textbf{Technique} \\
\midrule
Put all $720$ student numbers in a generator and drew $48$. & \blank{3.2cm} \\
Split the school by grade and drew students at random from every grade, in proportion. &
  \blank{3.2cm} \\
Ordered the list and took every $15$th student after a random start. & \blank{3.2cm} \\
Drew $2$ advisory groups at random and surveyed everyone in them. & \blank{3.2cm} \\
\end{tabularx}

\vspace{4pt}
(b) For a \textbf{systematic} sample of $48$ students, show why $k = 15$, then list the first
four students chosen from a random start of $7$.

\begin{work}
  k &= N \div n = 720 \div 48 = 15
\end{work}

\vspace{-2pt}
\blank{1.6cm} \quad \blank{1.6cm} \quad \blank{1.6cm} \quad \blank{1.6cm}

\vspace{4pt}
(c) For a \textbf{simple random} sample, a generator gives $471$, $839$, $260$, $554$, $128$.
One of those numbers cannot be used. Which one, and why?

\par\writeline

\end{enumerate}
\normalsize

% ======================================================================
\vspace{4pt}
\boxguard[10]
\parthead{Part D --- Extended Response \hfill 1 question, 10 points}

\small
\begin{enumerate}[leftmargin=2.2em, itemsep=10pt, topsep=2pt, start=19]

\item \textbf{Millbrook} ($720$ students) wants to know whether students would use a late
activity bus. The office has a numbered list of all $720$ students, the four grades are
believed to differ a lot in their answers, and the $30$ advisory groups are \textbf{built by
grade}.

\vspace{2pt}
(a) Choose \textbf{one} of the four sampling techniques and justify it from what this context
allows. Your reason must name \emph{how the groups are built}.

\par\writeline
\par\writeline

\vspace{2pt}
(b) Instead, a teacher draws $2$ advisory groups at random and surveys all $48$ students in
them; $30$ say yes. Show her two computations: the percent of her sample saying yes, and the
number of all $720$ students that percent points to.

\begin{work}
  \dfrac{30}{48} &= 0.625 = 62.5\% \\
  0.625 \times 720 &= 450 \text{ students}
\end{work}

(c) Her arithmetic is correct. Explain what is wrong with her report anyway, and state what
she \emph{is} allowed to say.

\par\writeline
\par\writeline
\par\writeline

\vspace{2pt}
(d) Is her $62.5\%$ a parameter or a statistic? Name one change to her method that would make
it an estimate worth trusting.

\par\writeline
\par\writeline

\end{enumerate}
\normalsize

\end{document}
